% !TeX root = ./Roadmap_instractions_main_v10.tex
% chktex-file 8 13 24 44

% =========================================================================
% REQUIRED PACKAGE ADDITIONS TO PREAMBLE (if not already present):
%   \usepackage{colortbl}   % for \rowcolor and \cellcolor
%   \usepackage{xcolor}     % for \definecolor and color names
%   \usepackage{array}      % for extended column formatting
%   \usepackage{booktabs}   % already present in v6
%   \usepackage{multirow}   % for \multirow in phase label column
% =========================================================================

% =========================================================================
% SECTION: Master Capability Transition Framework
% Insert before Section 6 (Detailed Priority Profiles)
% =========================================================================

\section{Methodology Overview}

\subsection{Master Capability Transition Framework - Transition Matrix}
\label{sec:master_transition}


\definecolor{gaterow}{gray}{0.91}
\definecolor{currentrow}{RGB}{253,243,237}
\definecolor{phase1row}{RGB}{237,245,253}
\definecolor{phase2row}{RGB}{237,249,244}
\definecolor{phase3row}{RGB}{245,244,254}
\definecolor{endstaterow}{RGB}{238,237,254}
\definecolor{sslred}{RGB}{251,235,235}
\definecolor{sslamber}{RGB}{250,238,218}
\definecolor{ssltrans}{RGB}{230,241,251}
\definecolor{sslteal}{RGB}{225,245,238}
\definecolor{sslpurple}{RGB}{238,237,254}

The ten executive priorities identified through the 4-Filter Prioritization
Framework in Section~\ref{sec:priorities} do not mature independently or at
uniform rates. Different priorities enter Phase~1 at distinct Operational
Deployment Phase (ODP) levels: Priorities~3, 6, 8, and~9 commence at the
Applied R\&D stage because their enabling technologies are not deployable
as of 2025--2026, whereas Priorities~1, 2, 4, 5, 7, and~10 commence at
Fragmented Piloting under regulatory activation. This heterogeneity is a
structural feature of the roadmap and must be preserved across all
per-priority analyses; collapsing it to a single Phase~1 characterisation
would constitute an overclaim of maturity for the less advanced priorities.

To provide a coherent methodology for tracking capability progression across
all ten priorities simultaneously, this section introduces a \textbf{Master
Capability Transition Framework} structured as a five-row, four-column matrix.
This framework is the primary analytical instrument of the roadmap and
constitutes the reference architecture against which all per-priority
progression matrices are instantiated. The matrix is not a summary of those
per-priority matrices: it defines the \textit{transition logic} governing
progression through the roadmap, including the explicit gate conditions that
must be satisfied before any priority can advance from one phase to the next.

\subsection*{Y-Axis: Phase States and Gate Conditions}

The Y-axis defines five states of capability maturity. The first state
(Current Status) establishes the 2025 baseline across all analytical
dimensions. Three sequential phase states follow, corresponding to the
roadmap's temporal structure: Phase~1 (2025--2030, Foundation and
Compliance), Phase~2 (2030--2035, Industrial Scaling, Standardisation and
Cross-Sector Deployment), and Phase~3 (2035--2040, Autonomous Leadership
and Frontier Innovation). The fifth state (Strategic End-State Objective,
2040 Vision) describes the target condition that the preceding three phases
collectively work to achieve.

Between each consecutive row pair, an explicit \textbf{gate condition} is
defined. A gate condition is the named, verifiable event --- a specific
regulatory enforcement date, a named standard finalization, or a measured
technical feasibility threshold --- whose activation is the necessary
condition for transition to the subsequent phase. Gate conditions are
derived exclusively from the Critical Enabler cells of the ten per-priority
progression matrices and are not aspirational descriptions. If a gate
condition has not been activated, the corresponding priority has not
transitioned, regardless of elapsed calendar time.

Each phase row is annotated with a \textbf{Strategic Sovereignty Level (SSL)}
band indicating the expected foreign-dependency trajectory during that phase.
The SSL band is a row-level property rather than a full column because
sovereignty evolves at a rate that is largely co-determined by the four
analytical dimensions simultaneously; representing it as a separate column
would create structural redundancy with the per-priority ODP/SSL matrices.
The SSL progression across the five rows is: High Foreign Dependency
(Current Status and early Phase~1) $\rightarrow$ Emerging EU Ecosystem
(late Phase~1 into Phase~2) $\rightarrow$ Full Sovereign Autonomy
(Phase~3) $\rightarrow$ Full Sovereign Autonomy with Export Leadership
(Strategic End-State).

\subsection*{X-Axis: Four Analytical Dimensions}

The X-axis defines four analytical dimensions capturing distinct transition
logics that operate simultaneously across each phase. \textbf{Research and
Innovation Maturity} tracks the scientific and technological advancement
trajectory, from academic fragmentation through applied and industrial R\&D
to frontier research, using TRL stage as a calibration anchor.
\textbf{Policy, Regulatory and Standards Enablement} tracks the activation
of the EU regulatory stack --- CRA, NIS2, AI Act, DORA --- and the maturation
of standardisation outputs from ETSI, CEN/CENELEC, 3GPP, and NIST, which
function as gate conditions for deployment transitions. \textbf{Operational
Deployment and Adoption} tracks observable deployment reality: what is
actually deployed, by which actor categories (national agencies, critical
infrastructure operators, research consortia, SMEs), and at what geographic
and sectoral scale. \textbf{Market Uptake and Strategic Impact} tracks EU
market share in critical security components, SME participation economics,
competitive positioning against non-EU hyperscalers, and the strategic
autonomy consequences of those market dynamics. This column is explicitly
not a duplicate of the Operational dimension: a capability can be
operationally deployed while remaining commercially inaccessible to SMEs or
strategically dependent on non-EU supply chains, as the case of Intel-TDX-based
confidential computing in Phase~1 illustrates.

\subsection*{Row V: Triple Accountability Framework}

The Strategic End-State Objective row is accountable to three distinct
frameworks simultaneously and must not be reduced to any single one of them.
First, \textbf{SNS JU SRIA alignment}: Priorities~6, 7, 9, and~10 must
satisfy the SNS JU SRIA Phase~3 security targets for 6G and the computing
continuum~\cite{snsju}. Second, \textbf{ECCC Strategic Agenda alignment}:
all ten priorities must contribute to the ECCC's mandate for European
industrial cybersecurity competitiveness and reduced foreign dependency,
measured through SSL metrics~\cite{eccc}. Third, \textbf{EU regulatory
stack coherence}: the full suite of NIS2, CRA, AI Act, DORA, Data Act, EU
Cloud Rulebook, and the envisioned ``Schengen for Cloud'' framework must
operate as a coherent, machine-enforceable system by 2040, with no
compliance gap between deployed sovereign capability and legal
obligation~\cite{nis2,cra,aiact}.

Table~\ref{tab:master_transition} presents the full Master Capability
Transition Framework. Key performance indicators cited in individual cells
are derived directly from the per-priority progression matrices in
Section~\ref{sec:priorities} and are not independently generated; they
serve as cross-reference anchors enabling consistency verification between
this master framework and the detailed priority profiles.

\newpage

\begingroup
\setlength{\LTleft}{0pt}
\setlength{\LTright}{0pt}
\scriptsize
\begin{longtable}{|p{1.7cm}|p{2.2cm}|p{2.2cm}|p{2.2cm}|p{2.2cm}|}
\caption{Master Capability Transition Framework: Phase States, Gate
Conditions, and SSL Bands across Four Analytical Dimensions}
\label{tab:master_transition} \\
\hline
\rowcolor{white}
\textbf{Phase / SSL Band} &
\textbf{R\&I Maturity} &
\textbf{Policy, Regulatory \& Standards} &
\textbf{Operational Deployment \& Adoption} &
\textbf{Market Uptake \& Strategic Impact} \\
\hline
\endfirsthead
\hline
\rowcolor{white}
\textbf{Phase / SSL Band} &
\textbf{R\&I Maturity} &
\textbf{Policy, Regulatory \& Standards} &
\textbf{Operational Deployment \& Adoption} &
\textbf{Market Uptake \& Strategic Impact} \\
\hline
\endhead

% ---- CURRENT STATUS ----
\rowcolor{currentrow}
\textbf{I. Current Status}

\smallskip
\textit{2025 baseline}

\smallskip
\colorbox{sslred}{\footnotesize\textbf{SSL: High Foreign Dependency}}
&
\footnotesize
\textbf{Fragmented, academically dominated.}
EU R\&I outputs concentrated in Horizon~2020 consortia; sovereign
alternatives at TRL~3--5. Heavy dependency on US-led outputs: NIST PQC
process, Intel TDX/AMD SEV-SNP TEE architectures, US/Chinese foundation
models for SecOps. No coordinated EU-wide applied R\&I programme for
sovereign secure silicon or edge-native confidential computing.
&
\footnotesize
\textbf{Regulatory stack adopted but not enforced.}
NIS2 enforcement initiated (Oct.~2024); CRA implementing technical acts
pending; AI Act adopted with high-risk provisions not yet operational;
EUCS candidate scheme non-binding; EUCC operational but limited uptake;
ENISA PQC migration guidance absent; NIS2 transposition fragmented
across Member States.
&
\footnotesize
\textbf{No continental sovereign security infrastructure.}
TEE deployment limited to foreign hyperscaler implementations; PQC
confined to experimental cloud pilots; Simpl pre-launch with no
operational security MIMs; federated SOC architecture at bilateral
proof-of-concept stage only; Island Mode not standardised; supply chain
attestation manual and opaque in Open~RAN ecosystems.
&
\footnotesize
\textbf{Near-zero EU market share in critical security components.}
EU domestic production of critical secure silicon near 0\%; US/Asian
hyperscalers dominate cloud security tooling and foundation models; no
EU-sovereign LLM for SecOps; SME compliance costs for CRA/NIS2
unaddressed by support schemes; RISC-V commercial security products
absent from public procurement frameworks.
\\
\hline

% ---- GATE I -> II ----
\rowcolor{gaterow}
\multicolumn{5}{|p{0.97\linewidth}|}{%
\footnotesize\textit{%
$\Rightarrow$~\textbf{Gate condition I\,$\to$\,II\,:}
\textbf{R\&I:} Horizon Europe calls for sovereign secure silicon and
lightweight PQC issued; ECCC work programmes activated.
\textbf{Policy:} CRA implementing technical acts in force; NIS2
Art.~21 national implementing guidance published by competent
authorities.
\textbf{Operational:} At least one national pilot per priority domain
initiated under regulatory pull.
\textbf{Market:} ECCC-co-funded SME compliance support scheme
operational; IPCEI-CIS initial deployments contracted.}}
\\
\hline

% ---- PHASE 1 ----
\rowcolor{phase1row}
\textbf{II. Phase 1}

\smallskip
\textit{2025--2030}

\smallskip
\textit{Pilot Validation \& Applied Maturity}

\smallskip
\colorbox{sslamber}{\footnotesize\textbf{SSL: High $\to$ Emerging EU}}
&
\footnotesize
\textbf{Applied R\&D (P3,P6,P8,P9); Fragmented Piloting (P1,P2,P4,P5,P7,P10).}
Horizon Europe calls targeting sovereign secure silicon, lightweight PQC
for 8/16-bit MCUs, Simpl security MIMs, and LLMSecOps. RISC-V secure
element prototypes targeting TRL~6--7 under EU Chips Act co-funding.
RF fingerprinting and CSI-based authentication validated on
EU-manufactured constrained hardware~\cite{cloud_edge_roadmap}.
ECCC-funded crypto-agility middleware research initiated per
FIPS~203/204/205 adoption~\cite{eccc}.

\smallskip
\textit{KPI: RISC-V secure element prototypes at TRL~6--7; lightweight
PQC benchmarked on ARM Cortex-M4 class hardware.}
&
\footnotesize
\textbf{Regulatory activation creates compliance market pull.}
CRA hardware security implementing rules published; NIS2 Art.~21
implementing guidance issued; ENISA PQC sector-specific migration
guidelines published by 2026, mapping FIPS~203/204/205 to NIS2
obligations~\cite{nis2}; EU AI Act Art.~9 compliance frameworks
operational for high-risk security AI~\cite{aiact}; EUCS certification
rollout mandating TEE workload protection~\cite{eucs}; Simpl-Open
official launch with initial security MIM definitions; SBOM templates
standardised and mandated in EU-funded deployments from 2026~\cite{cra}.

\smallskip
\textit{KPI: ENISA PQC guidelines published and adopted by 5+
Member States by 2026.}
&
\footnotesize
\textbf{Pilot deployments across priority domains under regulatory pull.}
3--5 RISC-V secure element pilots under Horizon Europe targeting EUCC
`High'; LLM-assisted SOC pilots in 5+ national cybersecurity centres
under AI Act Art.~9; CAMARA API profiles deployed in 3+ cross-border
5G corridor pilots by 2027; federated SIEMs with MITRE ATT\&CK
integration operational; Simpl security pilots in 2+ Common Data
Spaces (Manufacturing, Health); Island Mode fallback tested in
NIS2/DORA-compliant edge architectures~\cite{concordia}.

\smallskip
\textit{KPI: 2+ Common Data Spaces piloting Simpl security MIMs;
federated SIEMs active in 3+ Member States.}
&
\footnotesize
\textbf{Sovereignty trajectory established; SME access initiated.}
SME compliance support scheme co-funded by ECCC and national agencies,
providing subsidised tooling and certification assistance; IPCEI-CIS
initial deployments activating bilateral operator interconnection; EU
secure silicon market share below 10\% but industrial investment
trajectory established under CRA pull; EUCC `High' referenced in
public procurement frameworks for critical infrastructure ICT.

\smallskip
\textit{KPI: EU secure silicon market share $<$10\% but investment
trajectory established; SME support scheme operational in 15+
Member States.}
\\
\hline

% ---- GATE II -> III ----
\rowcolor{gaterow}
\multicolumn{5}{|p{0.97\linewidth}|}{%
\footnotesize\textit{%
$\Rightarrow$~\textbf{Gate condition II\,$\to$\,III\,:}
\textbf{R\&I:} RISC-V secure element prototypes validated at TRL~6--7;
lightweight PQC performance benchmarked on EU-manufactured constrained
hardware.
\textbf{Policy:} ENISA PQC guidelines published; EUCS certification
rollout active; Simpl-Open launched with security MIM definitions.
\textbf{Operational:} Federated SIEM pilots with standardised CTI
sharing operational in 3+ Member States; 2+ Common Data Space
Simpl pilots validated.
\textbf{Market:} CRA enforcement generating measurable industrial
investment in sovereign hardware alternatives.}}
\\
\hline

% ---- PHASE 2 ----
\rowcolor{phase2row}
\textbf{III. Phase 2}

\smallskip
\textit{2030--2035}

\smallskip
\textit{Industrial Scaling, Standardisation \& Cross-Sector Deployment}

\smallskip
\colorbox{ssltrans}{\footnotesize\textbf{SSL: Emerging EU $\to$ Full Autonomy}}
&
\footnotesize
\textbf{Applied and industrial R\&D concurrent with standardisation.}
Phase~2 is not purely a standardisation phase: hardware--software
co-design for neuromorphic/PIM architectures proceeds (P3); sovereign
EU LLM/SLM for SecOps reaches AI Act high-risk system compliance (P6);
RISC-V TEE designs enter volume production via Chips~JU (P1, P8); FHE
performance optimisation programmes initiated (P8); ZKP-based supply
chain attestation R\&D (P5); Cognitive Twin prototype development (P9).

\smallskip
\textit{KPI: At least one RISC-V TEE design in volume production by
2032; sovereign EU LLM/SLM certified at AI Act high-risk level.}
&
\footnotesize
\textbf{Standardisation milestones convert pilots into binding frameworks.}
ETSI TC CYBER hybrid PQC profiles published by 2031 (ML-KEM per
FIPS~203; ML-DSA per FIPS~204~\cite{fips203,fips204});
CEN/CENELEC JTC~13 Simpl security MIMs ratified by 2032; 3GPP SA3/ETSI
TC ITS standardised Island Mode fallback APIs published by 2031;
Runtime SBOM mandate for critical infrastructure supply chains by 2031;
classical-only cryptography deprecated for high-security EU government
contexts by 2033; Data Visa legal instruments developed~\cite{snsju}.

\smallskip
\textit{KPI: CEN/CENELEC Simpl security MIMs ratified; ETSI hybrid
PQC profiles published; classical-only deprecated in sovereign
contexts by 2033.}
&
\footnotesize
\textbf{Continental-scale federated deployment across regulated sectors.}
First EU RISC-V secure element at EUCC `High' by 2031--2032~\cite{eucc};
hybrid PQC deployed across Industrial IoT, energy grids, and financial
systems; Pan-European CTI platform operational as LLMSecOps analytics
engine by 2032; UEBA baselines harmonised using STIX/TAXII and MITRE
ATT\&CK; multi-party privacy-preserving AI pipelines demonstrated in
health and finance sectors by 2033; AI-driven graceful degradation
standardised in O-RAN AI/ML framework; inter-operator east-west
security interfaces operational in 2+ cross-border corridors.

\smallskip
\textit{KPI: 70\%+ Common Data Spaces on Simpl-based security; APT
mean dwell time reduced by 70\% relative to Phase~1 baseline.}
&
\footnotesize
\textbf{EU industrial sovereignty consolidating; critical component share rising.}
EU Chips Act fabrication producing RISC-V TEE designs in volume; EU
secure silicon reaching 30--40\% market share in critical infrastructure
procurement; sovereign European edge LLM/SLM in operational SecOps
deployment; IPCEI-CIS outcomes translated into binding inter-operator
SLA frameworks with security audit requirements; Data Visa pilots
operational in Nordic--Baltic and BeNeLux corridors; competitive
NaaS offers from European consortia viable.

\smallskip
\textit{KPI: EU secure silicon 30--40\% critical infrastructure
share; 70\%+ Data Spaces on Simpl security interoperability.}
\\
\hline

% ---- GATE III -> IV ----
\rowcolor{gaterow}
\multicolumn{5}{|p{0.97\linewidth}|}{%
\footnotesize\textit{%
$\Rightarrow$~\textbf{Gate condition III\,$\to$\,IV\,:}
\textbf{R\&I:} First EU RISC-V TEE design certified at EUCC `High'
and in volume production; sovereign EU edge LLM/SLM at AI Act
high-risk compliance.
\textbf{Policy:} CEN/CENELEC Simpl security MIMs ratified; ETSI
hybrid PQC profiles published; 3GPP Island Mode APIs standardised.
\textbf{Operational:} Pan-European CTI platform operational as
LLMSecOps engine; EU secure silicon exceeding 30\% critical
infrastructure share.
\textbf{Market:} EU NaaS consortium offers commercially competitive
in at least two cross-border corridors.}}
\\
\hline

% ---- PHASE 3 ----
\rowcolor{phase3row}
\textbf{IV. Phase 3}

\smallskip
\textit{2035--2040}

\smallskip
\textit{Autonomous Leadership \& Frontier Innovation}

\smallskip
\colorbox{sslteal}{\footnotesize\textbf{SSL: Full Sovereign Autonomy}}
&
\footnotesize
\textbf{Frontier research defining global state of the art.}
FHE reaching operational edge latency thresholds; zero-energy air
interface security with ambient RF harvesting; self-evolving AI swarm
protocols for autonomous Moving Target Defense; Cognitive Twin E2E
orchestration for cyber-physical resilience; bio-inspired elasticity
protocols standardised via ETSI ENI; battery-less node quantum-resistant
authentication. EuroHPC hosts federated training infrastructure for
self-improving sovereign security AI, preventing regression to
non-EU foundation model dependency.

\smallskip
\textit{KPI: FHE latency within $2\times$ AES-GCM baseline for edge
analytics; EuroHPC sovereign security AI model in operational
deployment.}
&
\footnotesize
\textbf{EU transitions from standard-taker to global standard-setter.}
EU Cloud Rulebook binding with machine-readable compliance APIs by
2037--2038; EUCC `High' as default procurement requirement for all
new EU ICT products; PQC-by-default mandated for all new EU ICT
products by 2036; AI Act explainability requirements enforced for
autonomous MTD decisions; ``Schengen for Cloud'' legal framework
enacted, harmonising data jurisdiction across 27 Member States;
ETSI/CEN/CENELEC PQC standards promoted to ISO/ITU-T as global
reference implementations.

\smallskip
\textit{KPI: EU Cloud Rulebook binding by 2038; PQC-by-default
mandate in force by 2036; Schengen for Cloud enacted.}
&
\footnotesize
\textbf{Fully autonomous, continent-wide sovereign security operations.}
EU Chips Act sovereign fabrication meeting 80\%+ of European secure
silicon demand by 2039; FHE operationalised for real-time edge analytics;
end-to-end confidential continuum from battery-less IoT to HPC; Cognitive
Twins autonomously coordinating E2E cyber-physical recovery across
critical infrastructure domains; AI-driven MTD continuously shuffling
network resources, IP addresses, and software diversity without human
intervention; SMPC standardised for cross-border federated AI in regulated
sectors (health, finance, defence, justice).

\smallskip
\textit{KPI: EU secure silicon self-sufficiency $\geq$80\% by 2039;
autonomous MTD and Cognitive Twin operational continent-wide.}
&
\footnotesize
\textbf{EU competitive parity achieved; technology standards exported globally.}
European operators achieving competitive scale parity with US/Chinese
hyperscalers in cloud-edge NaaS; EU-designed security components
dominant in European market; EUCS/EUCC adopted as global reference
certification frameworks; ``Schengen for Cloud'' eliminating national
fragmentation constraining European scale; EU cybersecurity technology
export balance turning positive; SME participation economically
viable without discriminatory compliance cost barriers.

\smallskip
\textit{KPI: EU NaaS competitive parity demonstrated in 3+ cross-border
corridors; EUCS/EUCC referenced in ISO/ITU-T standards.}
\\
\hline

% ---- GATE IV -> V ----
\rowcolor{gaterow}
\multicolumn{5}{|p{0.97\linewidth}|}{%
\footnotesize\textit{%
$\Rightarrow$~\textbf{Gate condition IV\,$\to$\,V\,:}
\textbf{R\&I:} FHE at operational edge latency; zero-energy
authentication validated on battery-less EU hardware; AI swarm
protocols stable and explainable under AI Act.
\textbf{Policy:} EU Cloud Rulebook binding; ``Schengen for Cloud''
enacted; PQC-by-default mandate in force; EU standards submitted
to ISO/ITU-T as reference implementations.
\textbf{Operational:} EU secure silicon self-sufficiency $\geq$80\%;
autonomous MTD and Cognitive Twin E2E orchestration operational
continent-wide.
\textbf{Market:} EU NaaS competitive parity demonstrated; technology
export balance positive.}}
\\
\hline

% ---- STRATEGIC END-STATE ----
\rowcolor{endstaterow}
\textbf{V. Strategic End-State Objective}

\smallskip
\textit{2040 Vision}

\smallskip
\colorbox{sslpurple}{\footnotesize\textbf{SSL: Full Autonomy + Export Leadership}}
&
\footnotesize
\textbf{Europe leads global security R\&I agenda.}
EU-originated R\&I outputs (PQC, TEE architectures, LLMSecOps,
zero-energy security, Cognitive Twins) define the global research
agenda and are adopted internationally. Alignment with SNS JU SRIA
Phase~3 outcomes: 6G security primitives sovereign and globally
deployed~\cite{snsju}. ECCC Strategic Agenda fully operationalised;
European industrial cybersecurity competitiveness self-sustaining
without emergency public subsidy~\cite{eccc}.

\smallskip
\textit{Triple accountability: SNS JU SRIA~\cite{snsju} $\cap$ ECCC
Strategic Agenda~\cite{eccc} $\cap$ EU regulatory stack~\cite{nis2}.}
&
\footnotesize
\textbf{Complete, coherent, machine-enforceable EU regulatory stack.}
NIS2, CRA, AI Act, DORA, Data Act, Cloud Rulebook, and ``Schengen
for Cloud'' operating as a coherent system with machine-readable
compliance APIs and no manual audit requirement. EUCS and EUCC
recognised globally as reference certification frameworks. ETSI and
CEN/CENELEC leading ISO/ITU-T security standardisation tracks.
Compliance is automated, verifiable in real time, and does not
impose discriminatory costs on SMEs or new market entrants.
&
\footnotesize
\textbf{Security-by-Design is the default operational baseline.}
Autonomous, self-healing, continent-wide security infrastructure
requiring no manual intervention for routine threat response. Full
confidential computing continuum from IoT sensor to HPC node. Zero
critical operational dependencies on non-EU security infrastructure.
Island Mode and graceful degradation are normative operating
assumptions, not emergency exceptions. Bio-inspired swarm-coordinated
resilience is the standard architecture for EU critical infrastructure.
&
\footnotesize
\textbf{Europe: sovereign, competitive, and globally influential.}
Europe is the global reference for sovereign, trustworthy, and
sustainable computing security. EU-designed security components
dominant in European market and competitive globally. SME participation
in EU cybersecurity ecosystem fully economically viable. European
security standards and AI-driven technologies exported as global
reference implementations. Strategic autonomy measured, sustained,
and no longer requiring emergency policy intervention to maintain.
\\
\hline

\end{longtable}
\endgroup

\subsection*{Phase~1 Heterogeneity: Applied R\&D vs. Fragmented Piloting}

A critical methodological distinction governs the Phase~1 row across the
ten executive priorities. Priorities~3 (Lightweight Cryptography and
Physical-Layer Security), 6 (Cognitive SecOps), 8 (Confidential Computing
and Privacy-Preserving AI Pipelines), and~9 (Resilient-by-Design
Infrastructure) commence Phase~1 in an Applied R\&D state because their
underlying enabling technologies --- zero-energy secure authentication,
EU-sovereign edge LLMs, RISC-V TEE hardware, and Cognitive Twin
orchestration respectively --- are not deployable at operational scale
as of 2025--2026. The remaining six priorities commence Phase~1 in a
Fragmented Piloting state because near-TRL solutions exist and regulatory
pull from CRA, NIS2, DORA, and the AI Act is sufficient to drive pilot
deployments in the near term.

This distinction is not cosmetic. Treating all ten priorities as uniform
pilots in Phase~1 would constitute a material overclaim of maturity to
both IEEE peer reviewers and European Commission technical evaluators, and
would produce internally inconsistent recommendations by implying that
technologies currently at TRL~3--4 can be mandated in operational
deployments by 2026. When instantiating per-priority progression matrices,
the R\&I cell for Priorities~3, 6, 8, and~9 in Phase~1 must state Applied
R\&D explicitly; it must not claim pilot or deployment activity.

\subsection*{Usage as a Consistency Reference}

The Master Capability Transition Framework serves three functions in the
broader roadmap document. First, it provides the structural skeleton within
which each per-priority progression matrix in Section~\ref{sec:priorities}
is located: every cell in a per-priority matrix should be traceable to the
corresponding intersection of phase row and analytical dimension in
Table~\ref{tab:master_transition}. Second, it provides a cross-priority
consistency check: if two per-priority matrices claim the same gate
condition (e.g., EU Chips Act fabrication capacity, which gates both
Priority~1 Phase~2 and Priority~8 Phase~2), that shared dependency should
be explicitly visible in the master framework and flagged in the executive
summary as a single-point-of-failure risk for the roadmap. Third, the
KPI anchors embedded in each phase cell provide a minimum measurable
outcome that must be demonstrated before the corresponding gate condition
can be declared satisfied; these KPIs are derived from the per-priority
matrices and are reproduced here as cross-reference anchors only.